\chapter{Serii de numere reale}

\section{Generalități}

O serie, înțeleasă informal ca o sumă infinită, este definită formal de două elemente:
\begin{itemize} \index{serii!șirul sumelor parțiale} \index{serii!șirul termenilor generali}
\item șirul termenilor generali;
\item șirul sumelor parțiale.
\end{itemize}

Astfel, de exemplu, dacă luăm seria $ \ds\sum_{n \geq 0} \frac{2^n}{n!} $, avem:
\begin{itemize}
\item șirul termenilor generali este $ x_n = \dfrac{2^n}{n!} $;
\item șirul sumelor parțiale este $ s_p = \ds\sum_{k = 0}^p \frac{2^k}{k!} $.
\end{itemize}

Seria se numește \emph{convergentă} dacă șirul sumelor parțiale este convergent,
iar limita acestui șir se numește \emph{suma seriei}. În caz contrar, seria
se numește \emph{divergentă}, adică șirul sumelor parțiale nu are limită sau
aceasta este infinită. Cînd vorbim despre \emph{natura seriei}, ne referim
dacă aceasta este convergentă sau divergentă.
\index{serii!convergente} \index{serii!divergente}
\index{serii!suma seriei}

Spre deosebire de cazul șirurilor din liceu, în cazul seriilor trebuie
mai multă atenție cînd discutăm convergența. Aceasta deoarece seriile au
o \emph{natură cumulativă}, adică toți termenii anteriori din șirul sumelor
parțiale își aduc contribuția. Mai precis, avem o recurență de forma:
\[
  s_{p + 1} = s_p + x_{p + 1}.
\]

De aceea, următoarele proprietăți sînt specifice seriilor:
\begin{enumerate}[(a)]
\item Dacă într-o serie schimbăm ordinea unui număr finit de termeni, obținem
  o serie nouă, care are aceeași natură cu cea inițială. Dacă există, suma seriei
  nu se schimbă.
\item Dacă eliminăm un număr finit de termeni dintr-o serie, se obține o
  serie nouă, cu aceeași natură. Dacă există, suma seriei poate să se schimbe.
\item Dacă o serie este convergentă, atunci ea are șirul sumelor parțiale mărginit.
\item Dacă o serie este convergentă, atunci șirul termenilor săi generali
  tinde către zero. Reciproca este, în general, falsă (contraexemplu $ \sum \frac{1}{n} $).

  Această proprietate ne permite să formulăm o \emph{condiție necesară de convergență}:
\item Dacă șirul termenilor generali ai unei serii nu este convergent către
  zero, atunci seria este divergentă.
\end{enumerate}
\index{serii!criteriu!necesar}

%%%%%%%%%%%%%%%%%%%%%%%%%%%%%%%%%%%%%%%%%%%%%%%%%%%%%%%%%%%%%%%%%%%%%%

\section{Serii cu termeni pozitivi}

În cazul seriilor care au șirul termenilor generali alcătuit numai
din numere pozitive, avem următoarele proprietăți, dintre care
unele rezultă prin particularizarea celor de mai sus:

\begin{enumerate}[(a)]
\item Șirul sumelor parțiale al unei serii cu termeni pozitivi este
  strict crescător.
\item O serie cu termeni pozitivi are întotdeauna sumă (finită sau nu).
\item O serie cu termeni pozitivi este convergentă dacă și numai dacă
  șirul sumelor parțiale este mărginit.
\item \textbf{Criteriul de comparație termen cu termen:} O serie care are termeni
  mai mari (doi cîte doi) decît o serie divergentă este divergentă.
  O serie care are termeni mai mici (doi cîte doi) decît o serie convergentă
  este convergentă. \index{serii!criteriu!comparație termen cu termen}
\item \textbf{Criteriul de comparație la limită, termen cu termen:}
  Fie $ \sum_n x_n $ și
  $ \sum_n y_n $ două serii cu termeni pozitivi. Presupunem că
  $ \frac{x_{n+1}}{x_n} \leq \dfrac{y_{n+1}}{y_n} $. Atunci:
  \begin{itemize}
  \item Dacă seria $ \sum_n y_n $ este convergentă, atunci și seria
    $ \sum x_n $ este convergentă;
  \item Dacă seria $ \sum_n x_n $ este divergentă, atunci și seria
    $ \sum_n y_n $ este divergentă.
  \end{itemize}
  \index{serii!criteriu!comparație la limită termen cu termen}
\item \textbf{Criteriul de comparație la limită:} Fie $ \sum_n x_n $ și
  $ \sum_n y_n $ două serii cu termeni pozitivi, astfel încît
  $ \ds\lim_{n \to \infty} \dfrac{x_n}{y_n} = \ell $.
  \begin{itemize}
  \item Dacă $ 0 < \ell < \infty $, atunci cele două serii au aceeași natură;
  \item Dacă $ \ell = 0 $, iar seria $ \sum_n y_n $ este convergentă, atunci și
    seria $ \sum_n x_n $ este convergentă;
  \item Dacă $ \ell = \infty $, iar seria $ \sum_n y_n $ este divergentă, atunci
    și seria $ \sum_n x_n $ este divergentă.
  \end{itemize}
  \index{serii!criteriu!comparație la limită}
\item \textbf{Criteriul radical:} Fie $ \sum_n x_n $ o serie cu
  termeni pozitiv, astfel încît $ \ds\lim_{n \to \infty} \sqrt[n]{x_n} = \ell $.
  Atunci:
  \begin{itemize}
  \item Dacă $ \ell < 1 $, seria este convergentă;
  \item Dacă $ \ell > 1 $, seria este divergentă;
  \item Dacă $ \ell = 1 $, criteriul nu decide.
  \end{itemize}
  \index{serii!criteriu!radical}
\item \textbf{Criteriul raportului:} Fie $ \sum_n x_n $ o serie cu termeni
  pozitivi și fie $ \ell = \ds\lim_{n \to \infty} \dfrac{x_{n+1}}{x_n} $.
  Atunci:
  \begin{itemize}
  \item Dacă $ \ell < 1 $, seria este convergentă;
  \item Dacă $ \ell > 1 $, seria este divergentă;
  \item Dacă $ \ell = 1 $, criteriul nu decide.
  \end{itemize}
  \index{serii!criteriu!raport}
\item \textbf{Criteriul lui Raabe-Duhamel:} Fie $ \sum_n x_n $ o serie
  cu termeni pozitivi și fie:
  \[
    \ell = \lim_{n \to \infty} n \cdot \Big( \frac{x_n}{x_{n+1}} - 1 \Big).
  \]
  Atunci:
  \begin{itemize}
  \item Dacă $ \ell < 1 $, seria este divergentă;
  \item Dacă $ \ell > 1 $, seria este convergentă;
  \item Dacă $ \ell = 1 $, criteriul nu decide.
  \end{itemize}
  \index{serii!criteriu!Raabe-Duhamel}
\item \textbf{Criteriul logaritmic:} Fie $ \sum_n x_n $ o serie cu termeni
  pozitivi și presupunem că există limita:
  \[
    \ell = \lim_{n \to \infty} \frac{\ln \frac{1}{x_n}}{\ln n}.
  \]
  Atunci:
  \begin{itemize}
  \item Dacă $ \ell < 1 $, seria este divergentă;
  \item Dacă $ \ell > 1 $, seria este convergentă;
  \item Dacă $ \ell = 1 $, criteriul nu decide.
  \end{itemize}
  \index{serii!criteriu!logaritmic}
\item \textbf{Criteriul integral:} Fie $ f: (0, \infty) \to [0, \infty) $ o
  funcție crescătoare și șirul $ a_n = \ds\int_1^n f(t) dt $. Atunci seria
  $ \sum_n f(n) $ este convergentă dacă și numai dacă șirul $ (a_n) $ este
  convergent.
  \index{serii!criteriu!integral}
\item \textbf{Criteriul condensării:} Fie $ (x_n) $ un șir astfel încît
  $ x_n \geq x_{n+1} \geq 0, \forall n $.
  Atunci seriile $ \sum_n x_n $ și $ \sum_n 2^n x_{2^n} $ au aceeași natură.
\end{enumerate}

%%%%%%%%%%%%%%%%%%%%%%%%%%%%%%%%%%%%%%%%%%%%%%%%%%%%%%%%%%%%%%%%%%%%%%

\section{Seria geometrică și seria armonică}

Două serii foarte importante pe care le putem folosi în comparații sînt
următoarele.

\textbf{Seria geometrică:} Fie $ a \in \RR $ și $ q \in \RR $. Considerăm
progresia geometrică de prim termen $ a $ și rație $ q $, care definește
seria $ \sum_n aq^n $, pe care o numim \emph{seria geometrică} de rație $ q $.

Suma parțială de rang $ n $ se poate calcula cu formula cunoscută din liceu:
\[
  S_n = a + aq + \dots + aq^{n-1} = %
  \begin{cases}
    a \cdot \dfrac{1 - q^n}{1 - q}, & q \neq 1 \\
    na, & q = 1
  \end{cases}
\]

Pentru convergență, să remarcăm că dacă $ |q| < 1 $, atunci:
\[
  \lim_{n \to \infty} s_n = \frac{a}{1 - q},
\]
deci seria este convergentă și are suma $ \dfrac{a}{1 - q} $.

Dacă $ |q| \geq 1 $, se poate verifica ușor, folosind criteriul
necesar, că seria geometrică este divergentă.
\index{serii!seria geometrică}

Un alt exemplu important este \textbf{seria armonică generalizată (Riemann)},
definită prin $ \ds\sum_n \dfrac{1}{n^\alpha} $, pentru $ \alpha \ in \RR $.
Se poate observa că:
\begin{itemize}
\item Dacă $ \alpha \leq 0 $, termenul general al seriei nu converge către
  zero, deci seria este divergentă, conform criteriului necesar;
\item Dacă $ \alpha > 0 $, termenii seriei formează un șir descrescător de
  numere pozitive. Seria este convergentă pentru $ \alpha > 1 $ și
  divergentă pentru $ \alpha \leq 1 $.
\end{itemize}

În cazul particular $ \alpha = 1 $, seria se numește simplu \emph{seria armonică}.
\index{serii!seria armonică}

%%%%%%%%%%%%%%%%%%%%%%%%%%%%%%%%%%%%%%%%%%%%%%%%%%%%%%%%%%%%%%%%%%%%%%

\section{Exerciții}

Studiați natura următoarelor serii cu termeni pozitivi, definite de șirul
termenilor generali $ (x_n) $, cu:
\begin{multicols}{2}
  \begin{enumerate}[(a)]
  \item $ x_n = \Big( \dfrac{3n}{3n + 1} \Big)^n $ (D, necesar);
  \item $ x_n = \dfrac{1}{n!} $ (C, comparație);
  \item $ x_n = \dfrac{1}{n \sqrt{n+1}} $ (C, comparație);
  \item $ x_n = \dfrac{n!}{n^{2n}}$ (C, raport);
  \item $ x_n = n! \Big( \dfrac{a}{n} \Big)^n, a \in \RR $ (discuție, raport);
  \item $ x_n = \dfrac{(n!)^2}{(2n)!} $ (C, raport);
  \item $ x_n = \dfrac{n!}{(a + 1)(a+2) \cdots (a+n)}, a > -1 $ (dis\-cu\-ți\-e, Raabe);
  \item $ x_n = \dfrac{1 \cdot 3 \cdot 5 \cdots (2n - 1)}{2 \cdot 4 \cdot 6 \cdots 2n} $ (D, Raabe);
  \item $ x_n = \Big( 1 - \dfrac{3\ln n}{2n} \Big)^n $ (C, logaritmic);
  \item $ x_n = \Big( \dfrac{n + 1}{3n + 1} \Big)^n $ (C, rădăcină);
  \item $ x_n = \Big( \dfrac{1}{\ln n} \Big)^{\ln(\ln n)} $ (D, logaritmic);
  \item $ x_n = \dfrac{1}{n \ln n} $ (D, integral);
  \item $ x_n = \dfrac{1}{n \ln^2 n} $ (C, integral);
  \item $ x_n = \dfrac{\sqrt[n]{(n+1)(n+2) \cdots (n + n)}}{n} $ (D, necesar + Cauchy);
  \item $ x_n = \dfrac{1}{7^n + 3^n} $ (C, comparație);
  \item $ x_n = \dfrac{1}{n \sqrt[n]{n}} $ (D, comparație);
  \item $ x_n = \sqrt{n^4 + 2n + 1} - n^2 $ (D, comparație);
  \item $ x_n = n^2 e^{-\sqrt{n}} $ (C, logaritmic);
  \item $ x_n = \dfrac{\ln n}{n^2} $ (C, comparație)
  \end{enumerate}
\end{multicols}


%%% Local Variables:
%%% mode: latex
%%% TeX-master: "../bcd-18-19"
%%% End:
